\begin{lemma}
\label{editorial-source-lemma-artinian-finite-nr-max}
If $R$ is Artinian then $R$ has only finitely many maximal ideals.
\end{lemma}

\begin{proof}
If there were infinitely many maximal ideals, choose distinct
$\mathfrak m_1,\mathfrak m_2,\ldots$ and put
$J_n=\bigcap_{i=1}^n\mathfrak m_i$ for $n\geq1$.
For every $i\leq n$, distinctness and maximality imply
$\mathfrak m_i+\mathfrak m_{n+1}=R$. Choose actual elements
$a_i\in\mathfrak m_i$ and $b_i\in\mathfrak m_{n+1}$ with
$a_i+b_i=1$. Their product
$a_1\cdots a_n$ belongs to every $\mathfrak m_i$ for
$i\leq n$, hence to $J_n$, and is congruent to $1$ modulo
$\mathfrak m_{n+1}$. It therefore does not belong to
$J_{n+1}=J_n\cap\mathfrak m_{n+1}$.
Consequently $J_1\supsetneq J_2\supsetneq\cdots$ is an
infinite strict descending chain of ideals, contradicting the
Artinian hypothesis. Thus the set of maximal ideals is finite.
For the zero ring the set is empty, which also satisfies the claim.
\end{proof}